% ============ 3. A Measurement-Decomposition Framework ============
\section{A Measurement-Decomposition Framework}\label{sec:framework}

The role of theory in this paper is narrow and deliberately so. We do not model price formation or claim to identify a mechanism (attention, liquidity provision, or otherwise). A natural first reading of the same-day co-movement of returns and sentiment is a price-formation story~---news moves prices~---but the audit below shows that the empirical signature feeding such interpretations, a large same-day loading of returns on sentiment, is substantially a property of the \emph{measured text}, not of the price process. What is needed instead is an accounting framework that separates the ingredients of the score so that each can be tested. This section provides it in three equations; every empirical section then maps one-to-one onto a term.

\subsection{The score equation}

Let $S^{N}_{i,t}$ denote the named-condition sentiment score for firm $i$'s headline at date $t$, and write
\begin{equation}\label{eq:score}
S^{N}_{i,t} \;=\; \underbrace{\kappa\,\theta_{i,t}}_{\text{semantic tone}}
\;+\; \underbrace{\iota\,G_i}_{\text{identity prior}}
\;+\; \underbrace{\rho\,P_{i,t}}_{\text{price echo}}
\;+\; \varepsilon_{i,t},
\end{equation}
where $\theta_{i,t}$ is the news-relevant fundamental signal in the text, $G_i$ is the scorer's firm-specific prior (any stable belief about firm $i$ that is independent of the headline at hand), $P_{i,t}$ is the realized price movement \emph{described by} the text (a retrospective headline such as ``rose 14\% in 30 days'' has $P_{i,t}\neq 0$ by construction), and $\varepsilon_{i,t}$ is call-level scoring noise. The anonymized score $S^{A}_{i,t}$ is generated by the same equation with the identity channel switched off ($\iota=0$ in effect, since the input no longer identifies the firm), and repeat calls identify $\Var(\varepsilon)$ directly. The dictionary score is a separate, low-duty-cycle instrument applied to the same text.

Equation~\eqref{eq:score} is an \emph{accounting} device, not a behavioral model: it imposes no orthogonality among the components. Indeed, the interesting case is the non-orthogonal one, in which news that matters for $\theta$ is also more likely to be retrospectively described and more likely to be covered at all~---the selection problem formalized in Section~\ref{sec:framework:selection}.

\subsection{What each audit test identifies}

The four scoring conditions deliver direct contrasts on the terms of \eqref{eq:score}:
\begin{align}
S^{N}_{i,t} - S^{A}_{i,t} &= \iota\,G_i + (\varepsilon^{N}_{i,t} - \varepsilon^{A}_{i,t})
&&\Rightarrow \text{H2: level shift and agreement (}\iota\text{),} \label{eq:contrast1}\\
S^{N}_{i,t}(r) - S^{N}_{i,t}(r') &= \varepsilon_{i,t}(r) - \varepsilon_{i,t}(r')
&&\Rightarrow \text{H3: repeat-call variance,} \label{eq:contrast2}
\end{align}
and the echo channel is tested by splitting the sample on a text-level flag rather than by a contrast of scores: if headlines that describe realized prices (the H4 flag) carry predictive loading that clean headlines lack, the loading is attributable to the $\rho P_{i,t}$ term of \eqref{eq:score}. The flag is a proxy for $P_{i,t}\neq 0$ at the headline level; it is not derived from \eqref{eq:score} and its validity is an empirical matter, which is why H4 is reported as a split-sample difference test rather than as an identity. The first contrast is contaminated by call-level noise, which is why the audit benchmarks the observed named-versus-anonymized dispersion against the repeat-call dispersion from the second contrast: only an excess of disagreement over the noise floor is evidence about $\iota$. This benchmarking is, to our knowledge, new to the LLM-finance setting. Two further contamination routes emerged from the staged designs, both measured rather than merely flagged. First, when the two conditions are fielded in different \emph{sittings}, the contrast also absorbs session-level drift $\eta$ that enters every score produced in a given sitting. The audit's own H2 estimate was initially confounded this way~---the named and anonymized rounds were twelve hours apart, and the drift alone exceeds the two-session shift (Section~\ref{sec:results:h2})~---which is why the same-session interleaving that repairs it is part of the protocol this paper ends up recommending. Second, the drift $\eta$ is itself two-thirds a \emph{batch-composition} effect~---the same headline scores differently depending on which other headlines share its batch~---rather than a property of the serving endpoint (Section~\ref{sec:results:h2}). The score of a headline is therefore a function of the text, the firm identity, the realized prices the text describes, the call, the sitting, and the batch: the audit measures each of the last three, and none of them is a property of the text.

For the return side, write the predictive regression of interest as
\begin{equation}\label{eq:pred}
r_{i,t+1} = \alpha + \beta\, \bar{S}_{i,t} + \gamma' X_{i,t} + u_{i,t},
\end{equation}
where $\bar{S}_{i,t}$ is the firm-day mean of the relevant condition's scores. Two consequences of \eqref{eq:score} discipline the interpretation of $\hat\beta$.

\begin{remark}[Attenuation and contamination, under explicit assumptions]\label{prop:attenuation}
Write $\bar{S}_{i,t}=\kappa\theta_{i,t}+\iota G_i+\rho P_{i,t}+\bar\varepsilon_{i,t}$ with $\bar\varepsilon$ independent of all other terms, and $r_{i,t+1}=b\,\theta_{i,t}+v_{i,t+1}$. Two restrictions make the statements below exact:
\begin{itemize}[leftmargin=1.8em, itemsep=1pt, topsep=2pt]
\item[\textbf{A1.}] $G_i$ is uncorrelated with $\theta_{i,t}$ and with $r_{i,t+1}$ conditional on the controls $X_{i,t}$.
\item[\textbf{A2.}] $P_{i,t}$ forecasts $r_{i,t+1}$ only through short-horizon continuation--reversal dynamics (a coefficient $b_P$, say), not through $\theta$.
\end{itemize}
A1 and A2 are \emph{testable restrictions, not maintained hypotheses}: H2 tests A1 and H4 tests A2, and Section~\ref{sec:results} reports the verdicts rather than assuming them. Under A1--A2, and additionally in the pure errors-in-variables case $G=P=0$: (i)~$\hat\beta \to b\kappa\,\lambda$ with $\lambda=\Var(\kappa\theta)/\Var(\bar{S})$ the reliability ratio~---the classical attenuation formula, which is why H3's ICC is a first-order input into $\hat\beta$'s magnitude. Outside that special case: (ii)~the price channel adds a contribution to the probability limit of $\hat\beta$ proportional to $\rho\,b_P\,\Var(P)/\Var(\bar{S})$~---exact only under mutual orthogonality of the components, which the framework deliberately does not impose, but directionally valid without it: an echo can generate predictive loading with no semantic content at all; and (iii)~the identity term shifts the probability limit of $\hat\beta$ only insofar as $G_i$ correlates with $r_{i,t+1}$ conditional on $X_{i,t}$~---exactly the object H2's predictive-difference test measures. We state these as a remark rather than as a proposition because the clauses are signposts with explicitly stated scope, not derivation output; clause~(ii) in particular must not be read as exact outside the orthogonality special case.
\end{remark}

The value of Remark~\ref{prop:attenuation} is signposting, not depth: clause (i)~says repeat-call reliability is a first-order input into interpreting $\hat\beta$'s magnitude (H3); clause (ii)~says the echo channel can generate a positive $\hat\beta$ with no semantic content at all, which is exactly what the H4 split tests; and clause (iii)~predicts the H2 outcome we observe~---no predictive shift in any of three contrasts of increasing design strength, with the level shift now measured at about three points by the interleaved design~---whenever firm priors are conditionally small in their return correlation. The same-day regression $r_{i,t}$ on $\bar S_{i,t}$ deserves its own remark: under \eqref{eq:score}, a large same-day loading is consistent both with prices moving on news \emph{and} with $P_{i,t}$ entering both sides. Observing that anonymization leaves the same-day loading intact (Section~\ref{sec:results:h4}) is what identifies the loading as a text property.

\subsection{Coverage selection}\label{sec:framework:selection}

The final channel operates before any scoring happens. Let $D_{i,t}\in\{0,1\}$ indicate that firm $i$ has a covered firm-day at $t$. If
\begin{equation}\label{eq:selection}
\Pr(D_{i,t}=1 \mid |r_{i,t}|,\ \sigma_i) \;=\; \Phi\!\left(\delta_0+\delta_1 |r_{i,t}|+\delta_2 \sigma_i\right),
\end{equation}
then the regression sample is drawn with intensity increasing in realized activity. The estimand of \eqref{eq:pred} is the conditional relationship \emph{within covered firm-days}; extrapolation to the full panel requires $\delta_1=\delta_2=0$, which H5 tests and rejects. We emphasize that this is a statement about \emph{what is estimated}, not an assertion that the estimate is ``wrong'': event-conditioned predictability is a legitimate object, but it is a different object from population predictability, and conflating the two is one of the quiet ways small-sample LLM studies overstate scope.

\subsection{Inference standard}\label{sec:framework:inference}

All regressions use two-way (firm and trading-day) clustered covariance following the short-panel convention; with 35 firms and 59 days, asymptotics are marginal, so H1 additionally reports a Rademacher wild-cluster bootstrap (day-level, null-imposed, $B=1{,}999$), instrument-difference tests use day-cluster bootstrap resampling ($B=999$), and the declared family (five hypotheses, seven primary comparisons) is reported with Benjamini--Hochberg false-discovery-rate $q$-values alongside raw $p$-values. The family itself was fixed after the first-round results were known but before any of the augmented inference above was run, and each later confirmatory stage was pre-specified in an archived design document before its own data collection. This is staged, \emph{within-study} pre-specification, not independent pre-registration, and we do not represent it as the latter. Readers should weigh two dangers together. On one side stands uncontrolled multiplicity: a pipeline like this one can generate dozens of contrasts, and undeclared selection among them would invalidate any $p$-value. On the other side stand the analyst degrees of freedom~---the garden of forking paths~---that any design chosen after seeing earlier results opens up: a staged study can be steered, one ``confirmatory'' step at a time, toward whatever the data reward, and pre-specifying each next stage in an archived document bounds but does not eliminate that risk. Our discipline is therefore threefold: report every contrast each stage produced, significant or not; keep the declared family separate from the staged additions and report both raw and BH-adjusted values for the family; and label exploratory statistics as exploratory. We adopt the convention~---standard in the econometrics literature, violated in this project's own first pass, and, we suspect, elsewhere besides~---that the difference between two instruments is never inferred from the separate significance of their coefficients.
